\section{Consistency checks at small sizes}\label{app:consistency}

The checks below compare the statements of the paper with exact computations of $T$ for small
$h$. They cannot confirm asymptotic statements, whose error terms are of the same size as the
differences observed, but they would have exposed errors in the finite forms of the bounds.
The polynomials $T$ were computed exactly by an implementation of the definitions in
\cref{sec:pencils} (interpolation of $\det[\mu_X(\varphi u^{i+j})]$ at $h+1$ rational points).

\subsection{The pole series}
\Cref{lem:pole-series} was tested on several entries $\mu_\xi(\varphi u^s)$ with the series
$\sum_mw^m(\varphi u^s)'(m)$ summed directly, with exact rational derivatives. For $\Gen$ at
$q=83$ and $\Band$ at $q=81$, $n=1$, both signs, the relative discrepancy is below $10^{-1240}$ at
$1500$ decimal digits (an earlier audit at $400$ digits found $10^{-60}$ and $10^{-145}$). For
$\Band$ the two parts of the affine entries are of size up to $10^{156}$ and cancel almost
completely at $X=\xi$; a first test at $120$ digits reported a spurious mismatch for this reason.

\subsection{The prime 2}
\Cref{tab:two-adic} compares $\vG2(T)$ with the finite bound of \cref{thm:prime-two} for $\Gen$ at
$q=83$ ($\vG2(T)$ is the same for both signs at these sizes).

\begin{table}[ht]
\centering
\caption{Exact $\vG2(T)$ for $\Gen$, $q=83$, against the finite bound of \cref{thm:prime-two}
and the leading term $-\frac{140}{99}h^2$. Last line: $\Band$, $q=81$, against
$-\frac{1847}{1100}h^2$.}\label{tab:two-adic}
\begin{tabular}{@{}rrrrr@{}}
\toprule
$h$ & $\vG2(T)$ & finite bound & leading term & $\vG2(T)/h^2$\\
\midrule
6 & $-40$ & $-43$ & $-50.9$ & $-1.111$\\
12 & $-172$ & $-180$ & $-203.6$ & $-1.194$\\
18 & $-405$ & $-417$ & $-458.2$ & $-1.250$\\
24 & $-730$ & $-748$ & $-814.5$ & $-1.267$\\
30 & $-1159$ & $-1177$ & $-1272.7$ & $-1.288$\\
36 & $-1696$ & $-1722$ & $-1832.7$ & $-1.309$\\
42 & $-2328$ & $-2358$ & $-2494.5$ & $-1.320$\\
48 & $-3051$ & $-3083$ & $-3258.2$ & $-1.324$\\
54 & $-4048$ & $-4080$ & $-4123.6$ & $-1.388$\\
60 & $-4816$ & $-4854$ & $-5090.9$ & $-1.338$\\
66 & $-5917$ & $-5966$ & $-6160.0$ & $-1.358$\\
72 & $-6998$ & $-7054$ & $-7330.9$ & $-1.350$\\
\midrule
40 ($\Band$) & $-2441$ & $-2462$ & $-2686.5$ & $-1.526$\\
\bottomrule
\end{tabular}
\end{table}

The ratios are not monotone, and all lie above $-140/99\approx-1.4141$. The entry at $h=54$ is an
outlier ($-1.388$, against $-1.324$ and $-1.338$ at $h=48,60$); it was recomputed on
26 November 2026 and gives the same value $-4048$. The finite bound shows the same dip
($-4080$), so it is a feature of the bound and not an error in the exact value; $h=54$ is the
size with $K=63=111111_2$, the largest binary digit sum of the $K$ in the table, which may explain the dip
through the low-order terms of \cref{thm:prime-two}; we have not isolated the mechanism.

For Krisman's determinant ($K=h$, $A=0$, $q=83$, $\sigma=1$) \cref{tab:two-adic-square} gives the
exact values in two normalisations: the un-normalised $\Delta=\det[\mu_X(\varphi u^{i+j})]$ and
$T=S\Delta$ with the normaliser $S$ of this paper, against the finite bound of
\cref{thm:prime-two} (valid for $T$).

\begin{table}[ht]
\centering
\caption{Exact $2$-adic Gauss valuations for $K=h$, $A=0$, $q=83$, $\sigma=1$; the limit of the
bound is $-\frac{13}{11}h^2$, that is $-1.1818h^2$. The columns with $/h^2$ show minus the
valuation divided by $h^2$.}\label{tab:two-adic-square}
\begin{tabular}{@{}rrrrrrr@{}}
\toprule
$h$ & $\vG2(\Delta)$ & $\vG2(T)$ & bound for $T$ & $-\vG2(\Delta)/h^2$ & $-\vG2(T)/h^2$ & $-$bound$/h^2$\\
\midrule
8 & $-66$ & $-42$ & $-48$ & $1.0312$ & $0.6562$ & $0.7500$\\
16 & $-287$ & $-223$ & $-236$ & $1.1211$ & $0.8711$ & $0.9219$\\
24 & $-663$ & $-583$ & $-604$ & $1.1510$ & $1.0122$ & $1.0486$\\
32 & $-1188$ & $-1028$ & $-1053$ & $1.1602$ & $1.0039$ & $1.0283$\\
40 & $-1882$ & $-1722$ & $-1760$ & $1.1763$ & $1.0762$ & $1.1000$\\
48 & $-2700$ & $-2492$ & $-2538$ & $1.1719$ & $1.0816$ & $1.1016$\\
56 & $-3685$ & $-3485$ & $-3533$ & $1.1751$ & $1.1113$ & $1.1266$\\
64 & $-4816$ & $-4432$ & $-4484$ & $1.1758$ & $1.0820$ & $1.0947$\\
\bottomrule
\end{tabular}
\end{table}

\subsection{Odd primes}
For $\Gen$ at $q=83$ and $h=6,12,\dots,42$, both signs, and for $\Band$ at $q=81$ and $h=40$, both
signs, every prime $p\le\max(4h,2K+2)+q$ was tested against the finite form of the relevant
bound: \cref{prop:p-divides-Y} for $p\mid q^2-4\sigma$, \cref{lem:large-primes} for $p>2K$,
\eqref{eq:vp-T-combined} with the exact class counts for $p^2>2K$, and \cref{prop:uniform}
otherwise. No bound failed. Typical margins: for $h=42$ the ordinary primes $p=41$ and $p=43$
have $\vG p(T)=-25$ and $-19$ against bounds $-26$ and $-20$; for $q^2-4=6885=3^4\cdot5\cdot17$ the
primes $3,5,17$ have $\vG p(T)=-1360,-422,-367$ against $-1560,-558,-390$. For $p=83$ one has
$\vG{83}(T)=h^2/3$ exactly for $h\le30$ and then slightly less ($0.3326h^2$ at $h=36$, $0.3206h^2$
at $h=54$, $0.3219h^2$ at $h=60$), as allowed by the term $-O(h\log h)$ in
\cref{prop:p-divides-q}. For $\Band$, $\vG3(T)=2106$ at $h=40$, against the leading term
$\frac43h^2=2133.3$.

\subsection{Nonvanishing}
For $\Gen$ at $q=83$, \cref{thm:nonvanishing} predicts $v_p(\Delta(0))=-18p+6$ for every
admissible prime $p$ with $h=6p$. Exact computation of $\Delta(0)$ gives $v_{17}(\Delta(0))=-300$
for $\sigma=-1$, $h=102$ (note $17\mid83^2-4$, so $p=17$ is excluded for $\sigma=1$), and
$v_{19}(\Delta(0))=-336$ for $\sigma=1$, $h=114$, as predicted. These are checks of the whole
mechanism of \cref{sec:nonvanishing}, including the reduction of the local functionals.

\subsection{Size}
The audit also computed $h^{-2}\log\abs{T(\xi)}$ in ball arithmetic. For $\Gen$ at $q=83$ it is
$-3.130$, $-3.111$, $-3.107$, $-3.107$, $-3.110$, $-3.111$, $-3.115$, $-3.116$ for
$h=12,\dots,54$ ($\sigma=1$), above the asymptotic bound $\bconst_{83}\le-3.1617$ by amounts
compatible with an error $O(\log h/h)$; for $\Band$ at $h=40$ it is $-3.377$ ($\sigma=1$) and
$-3.379$ ($\sigma=-1$), against $\bconst\le-3.3742$.

\subsection{The case \texorpdfstring{$q=79$}{q = 79}}
In floating point we recomputed the three numbers of \cref{sec:limitations}: the value
$-3.358974$ of $C+E(\mu_{79})+\int V\mu_{79}$, the trapezoidal bound $-1.985157$ for the profile
integral, and the resulting upper bounds $-3.365088$ and $-3.365333$ for the arithmetic
constant at $q=79$.
